\section{Decentering the Western Tradition}

I am among the last of a generation. Even those just a few years younger than I had a quite different experience. When a stone is thrown into a still pond, there is a wave motion that gradually attenuates, inversely to the square, the further it moves from the center, until the evidence of the original impulse is no longer recognizable. Attending public school systems in the Boston area, I was on one of those outer waves that was still under the influence of the first impulse. The mystery, then, to me is why so few of my generation have come to similar conclusions as mine.

Elementary school was in a plain brick building that could have been mistaken for a warehouse. Class sizes were around 25 to 30. A federal judge today would condemn it, as though erecting a new building in the modern style could make a difference to the education within. In elementary school, we began every day with the Pledge of Allegiance and the Our Father prayer, a tacit acknowledgment of spiritual authority and political power. One afternoon per week, we were let out of school and led to our respective nearby churches for an hour of religious education.

Yes, we did drills for a possible nuclear attack and were taught to hide under our desks and not look out the window to avoid being blinded. Obviously, if we were close enough to the bomb site to risk blindness, hiding under the desk was useless. Keep that in mind whenever the government implements a policy to keep the people safe. Beyond that, we were protected by Nike missiles and the Air Force was patrolling the skies; a sonic boom could occasionally be heard.

Our books were thick, without many pictures; unfortunately, I don't know that the book lists from that era are available. I recall we studied American and European history. We read the Odyssey in a text form. We read about the Germanic religion. We sang patriotic songs in music class. I don't know if the common core that is so popular now covers the same material.

There was no anti-German feeling, as the rocket engineers were well-respected. Even the rebels were respected and were frequently heroes of TV shows. Robert E Lee and Stonewall Jackson were well known names and considered to be fine Americans.

The school had a baseball field, where the boys played with no adult supervision. Even at the age of 8 or so, we were free to wander around the neighborhood, not expected home until the street lights came on. The kids had to provide their own entertainment. There was little point going home, since there was very little children's programming except for Saturday morning. In inclement weather, we played board and card games.

In junior high (middle) school, our sex education consisted in how to ask a girl for a date and what to say when meeting her parents. Since on the wave I mentioned, an educated man was expected to know ballroom dance and to speak French. So I was taught those in school. We were taught to speak proper grammar to get ahead in life, as well as habits such as diligence, punctuality, and so on. Clearly, that was another lie, since speaking and writing good English is wasted on the ignorant; as for the habits, the actual situation is slightly different, but that is a different topic.

We also had an intellectually rich home environment. My mother used to buy me books from the series Tom Corbett and the Space Cadets. From them, I learned about human types: the leader, the more intellectual navigator, and the jovial, but dimmer, engineer. In one episode, the computer went down, so the navigator had to compute the logarithms by hand. That appealed to me, because my father's table of logarithms and trig functions kept me entertained for hours. I couldn't get through his calculus book until I learned some algebra in the 7\textsuperscript{th} grade. I made it through differential calculus and integral calculus.

At home, I was fortunate that my parents had a collection of classic novels. I was able to read Edgar Allen Poe, Robert Lewis Stevenson, Herman Melville, James Fenimore Cooper, Nathaniel Hawthorne, Mark Twain, Daniel Defoe, Rudyard Kipling, and perhaps some others. Of course, there were the ones I didn't read: Louisa May Alcott, the Bronte's, Jane Austen. It was implicitly understood that they were for my sister to read.

\paragraph{The Expansion of Tradition}


So, that is what I envision when I hear about ``saving the West'', ``restoring Tradition'', etc., although I'm not sure how many others mean that. Obviously, the dominant ideology, not in a majoritarian sense but as the intellectual influence, is totally opposed to all that. Now we have twerking instead of ballroom, street lingo instead of good English, any other language but French, and a persistent, unrelenting attack on the implicit worldview\footnote{\url{http://www.thecollegefix.com/post/14698/}} of those authors I read, and whose worldview I absorbed. This is understandable coming from there, but there is a fifth column that is attacking from within, often unawares.

Several years ago, I was in correspondence with a Midwestern woman, a devout Evangelical, very conservative, who seemed educated and intelligent. She found this question game, and one of them was to name one's favorite historical figures. Mine were all literary, probably pompous as I looked back at the answers, except that I have actually read them. Hers took me by surprise, because it included Corrie Ten Boom and Harriet Tubman. It is hard to explain why, since they were both brave and admirable women, and worthy of her admiration. Nevertheless, I felt, at the time, that that selection was decentered, i.e., while important to their respective interests, it represented something outside the main thrust of the tradition that I had assumed she and I were discussing.

As I had time to think on it, I realized that the selection represented either an expansion of a wedge. For my friend, those two women were devout Christians. For her, the Holocaust and the emancipation of the slaves represented the expansion of the Western tradition beyond its eurocentrism. For the anti-tradition, the religion of those women was irrelevant; rather, the two issues they brought up act as a wedge to attack the entire West at its foundation. This is why all far right movements fail, since they are forced to exist within that shadow. The equation is Pro-West = Pro-Holocaust and Pro-Slavery. Mantras and complaints will not refute it because perception is reality, and that is the perception. In some cases, it is actually the reality.


\flrightit{Posted on 2013-09-30 by Cologero}

\begin{center}* * *\end{center}

\begin{footnotesize}
\begin{sffamily}

\texttt{Michael on 2013-09-30 at 21:08 said:}

It is fascinating to see how the rate of change has increased. When I grew up there was still enough of the vestiges of the old civilization to make things very pleasant.

Also fascinating to see how events can accelerate change. I am thinking of 9/11 and how it seemed to beget an evangelical atheism.

\hfill

\end{sffamily}
\end{footnotesize}

